\chapter{Primerjani modeli}
\label{sec:modeli}

V nalogi se osredotočamo na primerjavo štirih družin modelov strojnega učenja za klasifikacijo čustvenih oznak iz podatkov sledilnika pogleda.
V tem poglavju razložimo, katere so te družine in na kratko opišemo vsak uporabljen model znotraj družine.
Grafovske modele opišemo podrobneje, saj predstavljajo metodološko jedro naloge, poleg tega pa so za podatke sledilnika pogleda manj standardiziran pristop kot uveljavljeni klasični modeli.

Vsi modeli uporabljajo iste izvorne podatke, ista 10-sekundna okna in iste delitve na učne, validacijske in testne množice.
Razlika je v predstavitvi vhodnih podatkov.
Negrafovski osnovni modeli prejmejo agregirano tabelarično predstavitev istega okna, zamrznjeni prednaučeni predstavitveni modeli dobijo časovne vrste izbranih signalov, grafovski modeli pa delujejo na grafu, ki predstavlja časovno okno meritev.

\section{Pregled primerjanih pristopov}
\label{subsec:pregled-primerjanih-pristopov}

V tabeli~\ref{tab:pregled-primerjanih-modelov} je prikazano, katere modele uporabljamo, kateri družini vsak model pripada, kakšne podatke sprejme na vhod in kakšno vlogo ima v naši primerjalni študiji.

\begin{table}[!htbp]
    \centering
    \caption{Pregled primerjanih modelov in njihovih vhodnih predstavitev.}
    \label{tab:pregled-primerjanih-modelov}
    \footnotesize
    \setlength{\aboverulesep}{0.25ex}
    \setlength{\belowrulesep}{0.4ex}
    \begin{tabularx}{\textwidth}{@{}>{\raggedright\arraybackslash}m{0.22\textwidth} >{\raggedright\arraybackslash}m{0.23\textwidth} >{\raggedright\arraybackslash}m{0.20\textwidth} >{\raggedright\arraybackslash}X@{}}
        \toprule
        \textbf{Družina modelov} & \textbf{Model} & \textbf{Vhod v model} & \textbf{Vloga v primerjavi} \\
        \midrule
        \multirow[c]{6}{=}{\raggedright izhodiščna klasifikatorja} & Naključni & \multirow[c]{6}{=}{\raggedright brez učenja iz značilk} & spodnja meja naključnega ugibanja \\
        \cmidrule(lr){2-2}\cmidrule(l){4-4}
        & Večinski & & spodnja meja napovedovanja najpogostejšega razreda \\
        \cmidrule{1-4}
        \multirow[c]{3.6}{=}{\raggedright klasični modeli strojnega učenja} & \ModelName{SVM} & \multirow[c]{3}{=}{\raggedright agregirane značilke okna} & \multirow[c]{3}{=}{\raggedright klasični model strojnega učenja} \\
        \cmidrule(lr){2-2}
        & \ModelName{LightGBM} & & \\
        \cmidrule(lr){2-2}
        & \ModelName{MLP} & & \\
        \cmidrule{1-4}
        \multirow[c]{6}{=}{\raggedright zamrznjeni prednaučeni predstavitveni modeli} & \ModelName{GazeMAE} & koordinati pogleda $(x, y)$ & prednaučena predstavitev za množico \emph{samo pogled} \\
        \cmidrule(l){2-4}
        & \ModelName{MOMENT} & časovne vrste izbranih signalov & prednaučena predstavitev za množice s signali zenic oziroma več signali \\
        \cmidrule{1-4}
        \multirow[c]{10.5}{=}{\raggedright grafovski modeli} & \ModelName{GCN} & homogeni graf okna & osnovni grafovski model \\
        \cmidrule(l){2-4}
        & \ModelName{HeteroGCN-mean} & \multirow[c]{3.5}{=}{\raggedright heterogeni graf okna} & ločena obravnava relacij s preprosto fuzijo \\
        \cmidrule(lr){2-2}\cmidrule(l){4-4}
        & \ModelName{HeteroGCN-MLP} & & ločena obravnava relacij z učljivo fuzijo \\
        \cmidrule(l){2-4}
        & \ModelName{HeteroGCN-MLP-w} & heterogeni graf okna in značilke povezav & najkompleksnejša stopnja arhitekturne lestvice \\
        \bottomrule
    \end{tabularx}
\end{table}


\section{Negrafovski osnovni modeli}
\label{subsec:negrafovski-osnovni-modeli}

\subsection{Izhodiščna klasifikatorja}
\label{subsec:izhodiscna-klasifikatorja}

Večinski in naključni klasifikator predstavljata spodnjo primerjalno mejo.
Večinski klasifikator vedno napove večinski razred učne množice, naključni klasifikator pa razrede napoveduje z enako verjetnostjo ne glede na distribucijo podatkov.
Pri binarni klasifikaciji to pomeni verjetnost $50~\%$ za vsak razred.
Če model ne preseže rezultatov teh dveh klasifikatorjev, težko trdimo, da se iz vhodnih podatkov nauči uporabnih vzorcev.

\subsection{Klasični modeli strojnega učenja}
\label{subsec:klasicni-modeli-strojnega-ucenja}

Za pravično primerjavo negrafovskim osnovnim modelom podamo agregirane značilke istih signalov, kot jih prejme GNN.
V primeru množice \emph{vsi signali} so to koordinati pogleda $(x,~y)$, velikosti leve in desne zenice, oddaljenost od sledilnika pogleda ter informacija o fiksacijah.
Dodamo tudi normaliziran čas meritve znotraj okna.
Za vsako vhodno značilko izračunamo opisne statistike na ravni okna: povprečje, standardni odklon, minimum, maksimum, razpon, mediano, prvi kvartil, tretji kvartil in medkvartilni razmik.
Pri fiksacijah dodamo še število fiksacij, delež meritev v oknu, ki pripadajo fiksacijam, vsoto trajanj, povprečno trajanje in največje trajanje fiksacije.
V primeru množice \emph{vsi signali} je eno okno v tabelarični predstavitvi opisano z 68 agregiranimi značilkami.
Takšna pretvorba časovnega signala v opisne statistike na ravni okna oziroma odziva je običajen pristop pri klasičnih modelih strojnega učenja za podatke pogleda~\cite{bozak2024feature,lim2020emotion,tarnowski2020eye}.

\subsubsection{\ModelName{SVM}}

\ModelName{SVM} uporabimo kot klasični model z RBF jedrom nad agregiranimi značilkami okna~\cite{cortes1995support}.
RBF jedro modelu omogoča nelinearno ločevanje razredov, zato je \ModelName{SVM} uporaben kot močan negrafovski osnovni model.
Ker je \ModelName{SVM} občutljiv na velikostni red vhodnih značilk, agregirane značilke pred učenjem standardiziramo.

\subsubsection{\ModelName{LightGBM}}

\ModelName{LightGBM} uporabimo kot predstavnika gradientno spodbujenih odločitvenih dreves~\cite{ke2017lightgbm}.
Model je primeren za tabelarične podatke, zato je naravna primerjava za agregirane značilke časovnih oken.
Za razliko od \ModelName{SVM} in \ModelName{MLP} standardizacije vhodnih značilk ne potrebuje, saj odločitvena drevesa delujejo na pragovih posameznih značilk.

\subsubsection{\ModelName{MLP}}

Osnovni \ModelName{MLP} uporabimo kot primer osnovnega nevronskega modela~\cite{goodfellow2016deep}.
Vhod modela je vektor agregiranih značilk enega 10-sekundnega okna, izhod pa logiti za razrede ciljne naloge.
Arhitektura je namerno preprosta: dve skriti polno povezani plasti z aktivacijo ReLU in izpuščanjem, za njima pa izhodna linearna plast.
Model učimo z navzkrižno entropijo.

\section{Zamrznjeni prednaučeni predstavitveni modeli}
\label{subsec:zamrznjeni-prednauceni-modeli}

Poleg modelov nad ročno agregiranimi značilkami uporabimo tudi dva prednaučena kodirnika signalov: \ModelName{GazeMAE} za podatke pogleda in \ModelName{MOMENT} za ostale časovne vrste~\cite{bautista2021gazemae,goswami2024moment}.
Takšni modeli vhodnega okna ne pretvorijo v vnaprej določene statistike, temveč iz neposrednega signala izračunajo vektorsko predstavitev.
V tej nalogi ju uporabimo kot zamrznjena kodirnika, zato njunih parametrov med učenjem ne spreminjamo.
Na izračunanih predstavitvah nato učimo ločeno klasifikacijsko glavo \ModelName{MLP}.
Ta \ModelName{MLP} je arhitekturno drugačen, vendar ločen in nepovezan s samostojnim klasifikatorjem \ModelName{MLP} iz družine klasičnih modelov strojnega učenja.

\subsection{\ModelName{GazeMAE}}
\label{subsec:gazemae-model}

Za podatke pogleda uporabimo odprtokodni prednaučeni kodirnik \ModelName{GazeMAE}~\cite{bautista2021gazemae}.
Model prejme samo koordinati pogleda $(x,~y)$, zato ga samostojno uporabimo le za množico \emph{samo pogled}.
Ker je bil \ModelName{GazeMAE} učen na $2~s$ oknih, vzorčenih s frekvenco $500~Hz$, uporabljena $10~s$ okna najprej pretvorimo v obliko, ki jo model pričakuje.
Koordinati pogleda razdelimo na pet neprekrivajočih se $2~s$ kosov in ju prevzorčimo na ciljno frekvenco $500~Hz$.
Kodirnik temelji na časovno-konvolucijskih mrežah, ki ločeno kodirajo položaj in hitrost pogleda.
Iz predstavitev posameznih kosov nato z združevanjem pridobimo predstavitev celotnega okna s~512 razsežnostmi.

\subsection{\ModelName{MOMENT}}
\label{subsec:moment-model}

\ModelName{MOMENT} uporabimo kot splošni prednaučeni model za časovne vrste~\cite{goswami2024moment}.
V nasprotju z \ModelName{GazeMAE} ni vezan samo na koordinate pogleda, zato ga uporabimo za večkanalne časovne vrste signalov, ki pripadajo izbrani množici signalov v posameznem eksperimentu.
To so lahko velikosti zenic, oddaljenost od sledilnika pogleda in trajanje fiksacije.
Samostojno ga uporabimo le za množico \emph{samo zenici}.
Vhodno okno pred uporabo prevzorčimo na 512 časovnih korakov.
\ModelName{MOMENT} temelji na arhitekturi transformerjev za časovne vrste in v uporabljeni nastavitvi izračuna predstavitev okna s~768 razsežnostmi.

\subsection{Združene predstavitve \ModelName{GazeMAE} in \ModelName{MOMENT}}

Pri množicah \emph{pogled in zenici} ter \emph{vsi signali} uporabimo združeno predstavitev obeh kodirnikov.
Predstavitev \ModelName{GazeMAE} in predstavitev \ModelName{MOMENT} konkateniramo v en skupni vektor, ki je nato vhod v klasifikacijsko glavo \ModelName{MLP}.

V rezultatih zamrznjene prednaučene predstavitvene modele obravnavamo kot eno
družino modelov \ModelName{GazeMAE/MOMENT}.
Pri množici \emph{samo pogled} ta vrstica pomeni uporabo \ModelName{GazeMAE}, pri množici
\emph{samo zenici} uporabo \ModelName{MOMENT}, pri množicah \emph{pogled in zenici} ter
\emph{vsi signali} pa združeno predstavitev obeh kodirnikov. Za to smo se odločili zato, da ohranimo tabelo rezultatov pregledno.


\section{Grafovski modeli}
\label{subsec:grafovski-modeli}

Grafovski modeli uporabljajo grafovsko predstavitev okna iz poglavja~\ref{sec:grafovska-predstavitev}.
Pri isti množici signalov vse stopnje uporabljajo iste vozliščne značilke, iste relacije in ista časovna okna.
Razlikujejo se po tem, kako obravnavajo tipe povezav, kako združujejo relacijske predstavitve in ali uporabljajo značilke povezav za izračun naučenih uteži povezav.
Najprej definiramo osnovno GNN, nato pa s tremi nadgradnjami postopno gradimo končni model.
V tabeli~\ref{tab:arhitekturna-lestvica-gnn} so na kratko predstavljene razlike med vsemi štirimi stopnjami grafovskih modelov.

\begin{table}
    \centering
    \caption{Arhitekturna lestvica grafovskih modelov.}
    \label{tab:arhitekturna-lestvica-gnn}
    \small
    \begin{tabularx}{\textwidth}{@{}>{\raggedright\arraybackslash}p{0.24\textwidth} >{\centering\arraybackslash}p{0.10\textwidth} >{\centering\arraybackslash}p{0.12\textwidth} >{\centering\arraybackslash}p{0.10\textwidth} >{\raggedright\arraybackslash}X@{}}
        \toprule
        \textbf{Model} & \textbf{Ločene relacije} & \textbf{Fuzija relacij} & \textbf{Uteži povezav} & \textbf{Namen stopnje} \\
        \midrule
        \ModelName{GCN} & ne & -- & ne & preveri osnovno grafovsko predstavitev v homogenem grafu \\
        \ModelName{HeteroGCN-mean} & da & povprečje & ne & preveri vpliv ločene obravnave relacij \\
        \ModelName{HeteroGCN-MLP} & da & \ModelName{MLP} & ne & preveri vpliv naučene fuzije relacijskih predstavitev \\
        \ModelName{HeteroGCN-MLP-w} & da & \ModelName{MLP} & da & preveri vpliv naučenih predznačenih uteži povezav \\
        \bottomrule
    \end{tabularx}
\end{table}


\subsection{Skupni cevovod grafovskih modelov}
\label{subsec:skupni-deli-grafovske-arhitekture}

Grafovski modeli na vhod prejmejo graf glede na izbrano množico signalov, kot definirano v~\ref{sec:grafovska-predstavitev}.
Razsežnost vhodnih preslikav prilagodimo izbrani množici signalov -- ločeno za značilke vozlišč, pri modelu \ModelName{HeteroGCN-MLP-w} pa tudi za značilke povezav.

Vhodne vektorje najprej preslikamo v skupni latentni prostor z dvoslojnim \ModelName{MLP}.
Tako dobimo začetno predstavitev vsakega vozlišča, ki jo nato posodabljajo grafovske plasti.
Te plasti posredujejo sporočila po povezavah grafa in pri vsakem vozlišču združijo informacije iz njegove okolice.
V grafovskih blokih uporabimo aktivacijo GELU, rezidualne povezave, normalizacijo plasti in izpuščanje.
Ti gradniki niso predmet primerjave med grafovskimi modeli.
Primerjava se osredotoča na obravnavo relacij, fuzijo relacijskih predstavitev in uporabo značilk povezav.
Iz končnih predstavitev vozlišč nato s pozornostnim združevanjem\footnote{
Ker napovedujemo oznako celotnega okna, predstavitve vozlišč združimo v predstavitev grafa.
Uporabimo $\mathbf{h}_G = \sum_{v \in V} \alpha_v \mathbf{h}_v$, kjer je $\mathbf{h}_G$ predstavitev grafa, $\alpha_v$ utež pozornosti za vozlišče $v$, $\mathbf{h}_v$ pa predstavitev vozlišča $v$.
Uteži pozornosti izračunamo s softmax normalizacijo: $\alpha_v = \exp(g(\mathbf{h}_v)) / \sum_{u \in V} \exp(g(\mathbf{h}_u))$, kjer je $g$ naučena funkcija, ki predstavitvi vozlišča priredi skalarno oceno pomembnosti.
} izračunamo predstavitev celotnega grafa, to pa pošljemo skozi klasifikacijsko glavo \ModelName{MLP}.
Glava izračuna dva logita, enega za neprijetno in enega za prijetno valenco.
Napoved je razred z večjim logitom.
Pri poročanju verjetnosti logite normaliziramo s softmaxom.
Podrobnosti o številu plasti, velikosti skritih predstavitev, verjetnosti izpuščanja in ostalih hiperparametrih so opisane v dodatku~\ref{app:hiperparametri}.

Vse GNN treniramo od začetka do konca (ang. \emph{end-to-end}), kar pomeni, da se vsi učljivi deli modela optimizirajo hkrati glede na klasifikacijsko izgubo.
To vključuje začetno preslikavo značilk vozlišč, grafovske plasti, morebitno fuzijo relacij, pozornostno združevanje in klasifikacijsko glavo.
V nadaljevanju opišemo specifike vsake arhitekture v štiri-stopenjski arhitekturni lestvici, začenši z najenostavnejšim -- s homogenim \ModelName{GCN}.

\subsection{\ModelName{GCN}}
\label{subsec:model-gcn}

\ModelName{GCN} je v osnovi homogeni grafovski model.
Uporablja isto grafovsko predstavitev okna, kot je bila opisana v poglavju~\ref{sec:grafovska-predstavitev}, vendar ne ločuje časovnih, prostorskih in fiksacijskih relacij.
Vse povezave združi v en homogen graf, posredovanje sporočil pa izvede z običajnimi plastmi \texttt{GCNConv}.
Model ne uporablja značilk povezav in posledično naučenih uteži povezav.

Tehnično lahko model zapišemo kot:
\begin{equation}
    X
    \rightarrow
        MLP_{\text{vhod}}
    \rightarrow
    f_{\GCN}
    \rightarrow
    \text{READOUT}
    \rightarrow
    f_{\text{head}},
\end{equation}
kjer je $X$ matrika značilk vozlišč, $MLP_{\text{vhod}}$ vhodni dvoslojni perceptron, $f_{\GCN}$ zaporedje grafovskih konvolucijskih plasti, korak $READOUT$ združi predstavitve vozlišč v predstavitev celotnega grafa, $f_{\text{head}}$ pa iz te predstavitve izračuna napoved razreda.
TODO{Zamenjaj to enačbo z diagramom, ko ga imaš. Z diagramom lažje prikažeš že samo razliko med tem in naslednjim modelom in imaš potem majhen diagram v vsaki subsekciji in obkrožiš razlike pri vsaki nadgradnji}

\subsection{\ModelName{HeteroGCN-mean}}
\label{subsec:model-heterogcn-mean}

\ModelName{HeteroGCN-mean} dobimo tako, da \ModelName{GCN} nadgradimo s heterogeno grafovsko predstavitvijo.
Kot opisano v~\ref{sec:grafovska-predstavitev} tukaj ločimo različne tipe relacij.
Relacijske predstavitve vozlišč združimo v skupno predstavitev vozlišča s preprostim povprečenjem.
Namen te stopnje je izolirati vpliv ločevanja relacij.
Kot prikazano v diagramu~\ref{fig:diagram-heterogcn-mean} je torej sprememba glede na homogeni \ModelName{GCN} je zato samo v tem, da se sporočila najprej izračunajo ločeno po relacijah, nato pa združijo s povprečenjem.
\TODO{ko bom naredil diagram, ga dodam tu}

\subsection{\ModelName{HeteroGCN-MLP}}
\label{subsec:model-heterogcn-mlp}

\ModelName{HeteroGCN-MLP} dobimo tako, da \ModelName{HeteroGCN-mean} nadgradimo z naučeno fuzijo relacijskih predstavitev.
Različne predstavitve vozlišč -- tiste, ki so v danem modelu na voljo izmed časovne naprej, časovne nazaj, prostorske ter fiksacijske -- najprej konkateniramo, nato pa jih z dvoslojnim perceptronom združimo v skupno predstavitev vozlišča.
S tem dopuščamo, da model sam določi, kateri deli katere relacije so pomembnejši za končno predstavitev.
Kot prikazano v diagramu~\ref{fig:diagram-heterogcn-mlp} je torej sprememba  glede na \ModelName{HeteroGCN-mean} je torej v zamenjavi ročno določene fuzije z naučeno fuzijo relacijskih predstavitev.
\TODO{ko bom naredil diagram, ga dodam tu}

\subsection{\ModelName{HeteroGCN-MLP-w}}
\label{subsec:model-heterogcn-mlp-w}

\ModelName{HeteroGCN-MLP-w} dobimo tako, da \ModelName{HeteroGCN-MLP} nadgradimo z naučenimi utežmi povezav.
Za izračun uteži povezav ima model tri večslojne perceptrone: $f_w^{\text{čas}}$, $f_w^{\text{prostor}}$ in $f_w^{\text{fiksacija}}$.
Ti na podlagi vektorja značilk posamezne povezave, katerega komponente so povzete v tabeli~\ref{tab:znacilke-povezav}, izračunajo skalarno predznačeno utež, ki določa, kako močno posamezna povezava vpliva na posredovanje sporočil.
Pri posamezni množici signalov se uporabijo le funkcije za relacije, ki so v grafu prisotne.
Časovne povezave naprej in nazaj si delijo isto funkcijo $f_w^{\text{čas}}$, smer povezave pa je podana kot dodatna značilka povezave $\rho_{ij}$.
Formule za izračun in normalizacijo uteži so podane v dodatku~\ref{app:naucene-utezi-povezav}.
Namen te \emph{končne} stopnje je preveriti, ali modelu koristi razlikovati moč in predznak prispevka posameznih povezav znotraj relacije.
\TODO{ko bom naredil diagram, ga dodam tu}

\subsection{Izbira grafovske konvolucije}
\label{subsec:izbira-grafovske-konvolucije}

V vseh grafovskih modelih uporabimo plast \texttt{GCNConv}~\cite{kipf2017gcn}.
\texttt{GCNConv} je enostaven in uveljavljen operator posredovanja sporočil, ki ga lahko ohranimo nespremenjenega skozi celotno arhitekturno lestvico.
Alternativni operatorji, kot sta GraphSAGE ali GIN, bi v našem primeru odprli dodatne arhitekturne odločitve in s tem otežili neposredno primerjavo stopenj.

Izbiro smo preverili tudi z manjšim dodatnim poskusom na množici signalov \emph{pogled in zenici}.
Razlike med preizkušenimi operatorji so bile razmeroma majhne, \texttt{GCNConv} pa je dosegel primerljive in konsistentne rezultate tako pri osnovnem homogenem modelu kot tudi uteženi heterogeni arhitekturi.
Podrobnosti primerjave so podane v dodatku~\ref{app:izbira-grafovske-konvolucije}.

\section{Povzetek poglavja}
\label{subsec:povzetek-modelov}

V delu primerjamo štiri družine modelov različnih kompleksnosti in različnih arhitektur: izhodiščna klasifikatorja, klasične modele strojnega učenja, zamrznjena prednaučena kodirnika in grafovske modele.
Vhodne predstavitve za te modele posledično razdelimo na tri vrste: agregirane tabelarične značilke, neagregirane časovne vrste signalov in grafovske predstavitve podatkov sledilnika pogleda.
Grafovska lestvica štirih arhitektur preverja, kaj prispevajo sama grafovska struktura, ločena obravnava relacij, naučena fuzija relacijskih predstavitev in uteženost povezav.
V naslednjem poglavju opišemo evalvacijski protokol za primerjavo opisanih modelov na štirih množicah signalov.
